\subsection{DIRECT IMAGE AND INVERSE IMAGE OF A FILTER BASE}

Let \(\mathfrak{B}\) be a filter base on a set \(\mathbf{X}\) , and let \(f\) be a mapping of \(\mathbf{X}\) into a set \(\mathbf{X}'\) ; then \(f(\mathfrak{B})\) is a filter base on \(\mathbf{X}'\) , for the relation \(\mathbf{M} \neq \emptyset\) implies \(f(\mathbf{M}) \neq \emptyset\) , and we have \(f(\mathbf{M} \cap \mathbf{N}) \subset f(\mathbf{M}) \cap f(\mathbf{N})\) . If \(\mathfrak{B}_1\) is a base of a filter which is finer than the filter of base \(\mathfrak{B}\) , then \(f(\mathfrak{B}_1)\) is a base of a filter finer than the filter of base \(f(\mathfrak{B})\) (no. 3, Proposition 4).

PROPOSITION 10. If \(\mathfrak{B}\) is an ultrafilter base on a set X and if \(f\) is a mapping of X into a set \(\mathbf{X}'\) , then \(f(\mathfrak{B})\) is an ultrafilter base on \(\mathbf{X}'\) .

Let \(\mathbf{M}'\) be a subset of \(\mathbf{X}'\) . If \(\overline{f}^{1}(\mathbf{M}')\) contains a set \(\mathbf{M}\) of \(\mathfrak{B}\) , then \(\mathbf{M}'\) contains \(f(\mathbf{M})\) ; if not, then \(\complement \overline{f}^{1}(\mathbf{M}') = \overline{f}^{1}(\complement \mathbf{M}')\) contains a set \(\mathbf{N}\) of \(\mathfrak{B}\) (no. 4, Proposition 5) and therefore \(\complement \mathbf{M}'\) contains \(f(\mathbf{N})\) . Hence the result follows from Proposition 6 of no. 4.

Consider in particular the case where \(f\) is the canonical injection \(\mathbf{A} \to \mathbf{X}\) of a subset \(\mathbf{A}\) of a set \(\mathbf{X}\) . If \(\mathfrak{B}\) is a filter base on \(\mathbf{A}\) then \(f(\mathfrak{B})\) is a filter base on \(\mathbf{X}\) . The filter \(\mathfrak{F}\) on \(\mathbf{X}\) generated by \(f(\mathfrak{B})\) is called the filter generated by \(\mathfrak{B}\) when \(\mathfrak{B}\) is considered as a filter base on \(\mathbf{X}\) . If \(\mathfrak{B}\) is an ultrafilter base on \(\mathbf{A}\) it is also an ultrafilter base on \(\mathbf{X}\) by reason of Proposition 10.

Let us next examine whether the inverse image of a filter base is a filter base. Let \(\mathfrak{B}'\) be a filter base on a set \(\mathbf{X}'\) , and let \(f\) be a mapping of a set \(\mathbf{X}\) into \(\mathbf{X}'\) ; then \(\overline{f}^{1}(\mathfrak{B}')\) is a filter base on \(\mathbf{X}\) if and only if \(\overline{f}^{1}(\mathbf{M}') \neq \emptyset\) for each \(\mathbf{M}' \in \mathfrak{B}'\) . This is an immediate consequence of the relation \(\overline{f}^{1}(\mathbf{M}' \cap \mathbf{N}') = \overline{f}^{1}(\mathbf{M}') \cap \overline{f}^{1}(\mathbf{N}')\) and of Definition 3 of no. 3. This condition can also be expressed by saying that every set of \(\mathfrak{B}'\) meets \(f(\mathbf{X})\) {[}or that the trace of \(\mathfrak{B}'\) on \(f(\mathbf{X})\) is a filter base{]}. If this condition is satisfied, then \(f(\overline{f}^{1}(\mathfrak{B}'))\) is a base of a filter finer than the filter of base \(\mathfrak{B}\) .

If \(\mathfrak{B}\) is a filter base on \(\mathbf{X}\) it is clear that the above condition is satisfied by \(\mathfrak{B}' = f(\mathfrak{B})\) ; \(\overline{f}^1 (f(\mathfrak{B}))\) is then a base of a filter coarser than the filter of base \(\mathfrak{B}\) .

Let A be a subset of a set X, \(\varphi\) the canonical injection \(A \rightarrow X\) ; if B is a filter base on X then \(\overline{\varphi}^{1}(\mathfrak{B})\) is the same as \(B_{A}\) . If we express this as a filter base of A by means of the above condition, we recover part of Proposition 8 of no. 5.

